%=============================================================================
% ψ-Field Communications: Geometry-Locked, Spoof-Proof Information Transfer
% via Scalar Refractive Corridors
%
% Supporting research paper for Patent #8:
% "ψ-Field Communication and Information Systems"
% Inventor: Gary Thomas Alcock
%=============================================================================

\documentclass[twocolumn,superscriptaddress,prd,nofootinbib]{revtex4-2}

%--- Packages ----------------------------------------------------------------
\usepackage{amsmath,amssymb,amsthm}
\usepackage{mathtools}
\usepackage{bm}
\usepackage{physics}
\usepackage{graphicx}
\usepackage{xcolor}
\usepackage{hyperref}
\usepackage{cleveref}
\usepackage{booktabs}
\usepackage{multirow}
\usepackage{array}
\usepackage{siunitx}
\usepackage{tikz}
\usetikzlibrary{arrows.meta,positioning,decorations.pathmorphing,patterns,calc}
\usepackage{pgfplots}
\pgfplotsset{compat=1.18}
\usepackage{tcolorbox}
\tcbuselibrary{theorems,skins}
\usepackage{float}
\usepackage{enumitem}

%--- Custom commands ---------------------------------------------------------
\newcommand{\DFD}{\textsc{DFD}}
\newcommand{\psifield}{\psi}
\newcommand{\astar}{a_*}
\newcommand{\azero}{a_0}
\newcommand{\alphafine}{\alpha}
\newcommand{\ka}{\kappa_\alpha}
\newcommand{\Wfunc}{W}
\newcommand{\mufunction}{\mu}
\newcommand{\Phipot}{\Phi}
\newcommand{\avec}{\mathbf{a}}
\newcommand{\asq}{a^2}
\newcommand{\order}[1]{\mathcal{O}(#1)}
\newcommand{\bra}[1]{\langle #1 |}
\newcommand{\ket}[1]{| #1 \rangle}
\newcommand{\braket}[2]{\langle #1 | #2 \rangle}
\newcommand{\SNR}{\mathrm{SNR}}
\newcommand{\BER}{\mathrm{BER}}
\newcommand{\GLRT}{\mathrm{GLRT}}
\newcommand{\Prob}{\mathrm{P}}
\newcommand{\erfc}{\mathrm{erfc}}
\newcommand{\Corridor}{\mathcal{C}}
\newcommand{\Template}{\mathcal{T}}
\newcommand{\Fingerprint}{\mathcal{F}}

\newtheorem{theorem}{Theorem}
\newtheorem{lemma}[theorem]{Lemma}
\newtheorem{proposition}[theorem]{Proposition}
\newtheorem{corollary}[theorem]{Corollary}
\newtheorem{definition}{Definition}

%=============================================================================
\begin{document}

\title{$\psi$-Field Communications: Geometry-Locked, Spoof-Proof Information\\
Transfer via Scalar Refractive Corridors}

\author{Gary Thomas Alcock}
\affiliation{Independent Researcher}

\date{\today}

\begin{abstract}
We develop a complete communication theory for information transfer via
the scalar refractive field $\psi$ of Density Field Dynamics (DFD).
A $\psi$ corridor---a shaped region of refractive index
$n = e^\psi$ connecting transmitter and receiver---provides a
deterministic one-way phase velocity $c_1 = c\,e^{-\psi}$ and a
route-specific propagation delay
$\tau = (1/c)\int e^\psi\,ds$ that cannot be forged by an adversary
lacking access to the physical field geometry.
Information is encoded in a composite constellation mapping bits to
$(\delta\psi,\,\dot{\psi},\,\nabla\psi)$---the amplitude, phase rate,
and spatial gradient of the field perturbation.
We derive the channel capacity under thermal and quantum noise,
proving that geometry-locked authentication via corridor fingerprint
matching in a generalized likelihood ratio test (GLRT) achieves
exponentially vanishing false-acceptance probability with corridor
complexity.
Orthogonal corridor multiplexing is shown to support $N_c$ independent
channels with zero cross-talk in the gradient-orthogonal limit.
Comparisons with RF, free-space optical, quantum key distribution, and
neutrino communication demonstrate that $\psi$-comms occupy a unique
niche: medium-penetrating, synchronization-free, natively authenticated
links immune to electromagnetic jamming.
We present designs for underwater, space-to-ground, emergency mesh,
and financial timing backbone applications, and specify a tabletop
demonstrator using high-Q cavity $\psi$-pumping at the EM--$\psi$
coupling threshold $|\lambda - 1| \lesssim 3 \times 10^{-5}$.
\end{abstract}

\maketitle
\tableofcontents

%=============================================================================
\section{Introduction}\label{sec:intro}
%=============================================================================

\subsection{Motivation}

Every modern communication system---from 5G cellular to deep-space
telemetry---encodes information on electromagnetic carriers that
propagate through a medium characterized by dispersion, multipath
fading, and vulnerability to jamming and spoofing.
Timing is recovered from two-way exchange protocols susceptible to
man-in-the-middle attacks.
Addressing and routing are logical constructs overlaid on a physical
layer that provides no intrinsic authentication.
No existing communication modality binds addressing, timing, and
routing to a physically instantiated scalar geometry that a spoofer
cannot forge along an arbitrary path.

Density Field Dynamics (DFD)~\cite{Alcock2025DFD} introduces the
scalar refractive field $\psi(\mathbf{r},t)$ as the fundamental
gravitational degree of freedom on flat Euclidean space.
The optical metric
\begin{equation}
d\tilde{s}^2 = -\frac{c^2\,dt^2}{n^2} + d\mathbf{r}^2, \qquad
n(\mathbf{r},t) = e^{\psi(\mathbf{r},t)},
\label{eq:optical-metric}
\end{equation}
governs light propagation, with the one-way phase velocity
\begin{equation}
c_1(\mathbf{r}) = c\,e^{-\psi(\mathbf{r})}.
\label{eq:c1}
\end{equation}
A shaped $\psi$ corridor---a spatially structured perturbation
$\delta\psi(\mathbf{r},t)$ connecting two nodes---provides:
\begin{enumerate}[label=(\roman*)]
\item A deterministic, one-way propagation delay
$\tau = (1/c)\int_{\Corridor} e^\psi\,ds$,
\item A route-specific ``fingerprint'' in the gradient map
$\nabla\psi(\mathbf{r})$ along the path,
\item Native immunity to EM jamming, since the information carrier is
the scalar field, not the transverse EM field.
\end{enumerate}

This paper develops $\psi$-field communication as a complete
physical-layer technology: modulation formats, channel capacity,
security proofs, multiplexing architecture, and comparison with all
existing communication modalities.

\subsection{Relation to DFD Core Theory}

The DFD field equation in its general nonlinear form is
\begin{equation}
\nabla \cdot \bigl[\mu\!\bigl(|\nabla\psi|/\astar\bigr)\nabla\psi\bigr]
= -\frac{8\pi G}{c^2}(\rho - \bar{\rho}),
\label{eq:field-eq}
\end{equation}
where $\mu(x) = x/(1+x)$ is the uniquely derived crossover function,
and $\astar = 2a_0/c^2$ sets the gradient scale.
For communication applications, we operate in the \emph{high-gradient}
(Newtonian) regime $|\nabla\psi|/\astar \gg 1$ where $\mu \to 1$,
so the linearized form
\begin{equation}
\nabla^2 \delta\psi = -\frac{8\pi G}{c^2}\,\delta\rho
\label{eq:linearized}
\end{equation}
governs small perturbations about a background $\psi_0$.

The EM--$\psi$ coupling parameter $\lambda$ (Appendix R
of~\cite{Alcock2025DFD}) controls whether electromagnetic fields can
actively source $\psi$ oscillations.
Existing cavity stability constrains $|\lambda - 1| \lesssim 3 \times 10^{-5}$;
intentional optimization reaches $|\lambda-1| \sim 10^{-14}$.
Communication applications require only $|\lambda - 1| \neq 0$ at any
detectable level---the modulation depth of the $\psi$ perturbation
determines the signal-to-noise ratio, not the absolute coupling strength.

\subsection{Paper Outline}

Section~\ref{sec:channel-model} develops the $\psi$-corridor channel
model and propagation physics.
Section~\ref{sec:modulation} defines composite modulation constellations.
Section~\ref{sec:capacity} derives channel capacity under noise.
Section~\ref{sec:security} proves geometry-locked authentication security.
Section~\ref{sec:multiplexing} develops orthogonal corridor multiplexing.
Section~\ref{sec:comparison} compares with RF, optical, QKD, neutrino,
and gravitational-wave communication.
Section~\ref{sec:applications} presents application-specific designs.
Section~\ref{sec:demonstrator} specifies a laboratory demonstrator.
Section~\ref{sec:conclusions} summarizes.

%=============================================================================
\section{$\psi$-Corridor Channel Model}\label{sec:channel-model}
%=============================================================================

\subsection{Corridor Geometry and Field Structure}

\begin{definition}[$\psi$-Corridor]
A $\psi$-corridor $\Corridor_{AB}$ between nodes $A$ and $B$ is a
spatially structured perturbation $\delta\psi(\mathbf{r},t)$ satisfying:
\begin{enumerate}[label=(\alph*)]
\item $\delta\psi(\mathbf{r},t) \neq 0$ on a connected tubular region
$\mathcal{V}$ with $A,B \in \partial\mathcal{V}$,
\item The field equation~\eqref{eq:linearized} is satisfied within
$\mathcal{V}$ for a prescribed source distribution $\delta\rho(\mathbf{r},t)$,
\item The gradient map $\nabla\delta\psi(\mathbf{r})$ restricted to
$\mathcal{V}$ is uniquely determined by the source geometry.
\end{enumerate}
\end{definition}

The corridor has a natural parameterization by arc length $s$ along its
central axis, with cross-sectional profile described by the transverse
field distribution $\delta\psi_\perp(s, \boldsymbol{\rho})$ where
$\boldsymbol{\rho}$ denotes the transverse coordinates.

For a cylindrically symmetric corridor of radius $R$ and length $L$,
the lowest-order mode has the profile
\begin{equation}
\delta\psi(s, \rho) = \psi_0(s)\,J_0\!\left(\frac{j_{01}\rho}{R}\right)
\exp\!\left(-\frac{s^2}{2\sigma_s^2}\right),
\label{eq:corridor-profile}
\end{equation}
where $J_0$ is the zeroth Bessel function, $j_{01} \approx 2.405$,
$\psi_0(s)$ is the axial amplitude, and $\sigma_s$ characterizes the
longitudinal envelope.

\subsection{Propagation Delay}

The one-way propagation delay along corridor $\Corridor$ is
\begin{equation}
\tau_\Corridor = \frac{1}{c}\int_\Corridor e^{\psi(\mathbf{r})}\,ds.
\label{eq:delay}
\end{equation}
For small perturbations $|\delta\psi| \ll 1$ about background $\psi_0$:
\begin{equation}
\tau_\Corridor \approx \frac{e^{\psi_0}}{c}\left(L + \int_0^L
\delta\psi(s)\,ds\right)
= \tau_0 + \delta\tau,
\label{eq:delay-linearized}
\end{equation}
where $\tau_0 = Le^{\psi_0}/c$ is the background delay and
\begin{equation}
\delta\tau = \frac{e^{\psi_0}}{c}\int_0^L \delta\psi(s)\,ds
\label{eq:delay-perturbation}
\end{equation}
is the information-bearing delay perturbation.

\begin{proposition}[Delay Uniqueness]
\label{prop:delay-unique}
Two corridors $\Corridor_1$, $\Corridor_2$ with distinct gradient maps
$\nabla\delta\psi_1 \neq \nabla\delta\psi_2$ have distinct delay
profiles $\tau_{\Corridor_1}(s) \neq \tau_{\Corridor_2}(s)$ almost
everywhere along their paths.
\end{proposition}

\begin{proof}
If $\tau_{\Corridor_1}(s) = \tau_{\Corridor_2}(s)$ for all $s$, then
$\int_0^s [\delta\psi_1(s') - \delta\psi_2(s')]\,ds' = 0$ for all $s$,
implying $\delta\psi_1 = \delta\psi_2$ a.e., contradicting the distinct
gradient assumption.
\end{proof}

\subsection{Corridor Fingerprint}

\begin{definition}[Corridor Fingerprint]
The fingerprint of corridor $\Corridor$ is the vector-valued function
\begin{equation}
\Fingerprint_\Corridor(s) \equiv
\bigl(\delta\psi(s),\;\nabla_\perp\delta\psi(s),\;\partial_s\delta\psi(s)\bigr),
\label{eq:fingerprint}
\end{equation}
evaluated along the corridor axis.
\end{definition}

The fingerprint encodes the complete local field geometry at each point
along the path.
A receiver with access to high-precision $\psi$-gradiometry can measure
$\Fingerprint_\Corridor$ and compare it against a stored template for
authentication (Section~\ref{sec:security}).

\subsection{Noise Model}

The $\psi$ channel is subject to three noise sources:

\paragraph{Thermal noise.}
Thermal fluctuations in the local mass-energy distribution produce
stochastic $\psi$ perturbations.
For a detection bandwidth $B$ and detector temperature $T_d$:
\begin{equation}
\sigma_\psi^{\rm th} = \sqrt{\frac{k_B T_d}{M_\psi c^2}}\,\sqrt{B},
\label{eq:thermal-noise}
\end{equation}
where $M_\psi$ is the effective mass of the $\psi$-mode coupling to the
detector.

\paragraph{Quantum noise.}
The quantum vacuum fluctuation of $\psi$ sets a fundamental noise floor.
For a detector of volume $V_d$ and integration time $T$:
\begin{equation}
\sigma_\psi^{\rm qm} = \sqrt{\frac{\hbar G}{c^3 V_d}}\,\frac{1}{\sqrt{T}}.
\label{eq:quantum-noise}
\end{equation}
This is analogous to the gravitational-wave detector shot-noise limit
and represents the Heisenberg-limited sensitivity.

\paragraph{Gravitational background.}
Ambient gravitational signals (seismic, tidal, astrophysical) produce a
confusion background:
\begin{equation}
S_\psi^{\rm grav}(f) \approx \frac{4G}{c^4}\,S_h(f),
\label{eq:grav-background}
\end{equation}
where $S_h(f)$ is the strain spectral density of the gravitational-wave
background at frequency $f$.

The total noise spectral density is
\begin{equation}
S_\psi^{\rm tot}(f) = S_\psi^{\rm th}(f) + S_\psi^{\rm qm}(f) +
S_\psi^{\rm grav}(f).
\label{eq:total-noise}
\end{equation}

%=============================================================================
\section{Composite Modulation}\label{sec:modulation}
%=============================================================================

\subsection{Three-Dimensional Symbol Space}

The $\psi$-field provides three independent degrees of freedom for
encoding information at each symbol interval:
\begin{enumerate}[label=(\roman*)]
\item \textbf{Amplitude:} $\delta\psi(t)$ --- the local field
perturbation strength,
\item \textbf{Phase rate:} $\dot{\psi}(t) \equiv \partial_t\psi$ --- the
temporal derivative encoding phase chirps,
\item \textbf{Gradient:} $\nabla\psi(\mathbf{r},t)$ --- the spatial
derivative vector (3 components in 3D).
\end{enumerate}

\begin{definition}[Composite Constellation]
A $\psi$-field constellation $\mathcal{S}$ of order $M$ is a set of $M$
symbol vectors
\begin{equation}
\mathbf{s}_k = (\delta\psi_k,\;\dot{\psi}_k,\;\nabla\psi_k)
\in \mathbb{R}^5, \qquad k = 1,\ldots,M,
\label{eq:constellation}
\end{equation}
where $\nabla\psi_k \in \mathbb{R}^3$ denotes the gradient vector at
the receiver location.
Each symbol carries $\log_2 M$ bits.
\end{definition}

\subsection{Amplitude Modulation: $\delta\psi$-ASK}

The simplest modulation encodes bits in the amplitude of $\delta\psi$.
For $M$-level amplitude-shift keying ($\delta\psi$-ASK):
\begin{equation}
\delta\psi_k = \left(2k - 1 - M\right)\frac{\Delta\psi}{M-1}, \quad
k = 0,\ldots,M-1,
\label{eq:ASK}
\end{equation}
with inter-symbol spacing $\Delta\psi_{\rm min} = 2\Delta\psi/(M-1)$.
The bit-error rate for coherent detection in Gaussian noise of
variance $\sigma_\psi^2$ is
\begin{equation}
\BER_{\rm ASK} = \frac{2(M-1)}{M\log_2 M}\,Q\!\left(\frac{\Delta\psi_{\rm min}}{2\sigma_\psi}\right),
\label{eq:BER-ASK}
\end{equation}
where $Q(x) = \frac{1}{2}\erfc(x/\sqrt{2})$.

\subsection{Phase-Rate Modulation: $\dot{\psi}$-FSK}

Analogous to frequency-shift keying in RF, $\dot{\psi}$-FSK encodes
information in the temporal derivative:
\begin{equation}
\dot{\psi}_k(t) = \dot{\psi}_0 + k\,\Delta\dot{\psi},
\qquad k = 0,\ldots,M-1.
\label{eq:FSK}
\end{equation}
Orthogonality requires symbol duration $T_s$ satisfying
$\Delta\dot{\psi} \cdot T_s = 2\pi / T_s$, i.e.,
\begin{equation}
T_s = \sqrt{\frac{2\pi}{\Delta\dot{\psi}}}.
\label{eq:FSK-orthogonal}
\end{equation}

The detector measures $\dot{\psi}$ via differential $\psi$ sampling at
two closely spaced times $t$ and $t + \delta t$:
\begin{equation}
\dot{\psi}_{\rm meas} = \frac{\psi(t+\delta t) - \psi(t)}{\delta t}.
\label{eq:psi-dot-measurement}
\end{equation}

\subsection{Gradient-Key Modulation: $\nabla\psi$-QAM}

The spatial gradient $\nabla\psi = (\partial_x\psi, \partial_y\psi,
\partial_z\psi)$ provides a three-dimensional modulation space.
A gradient-QAM constellation places symbols on a 3D lattice in gradient
space:
\begin{equation}
\nabla\psi_k = \Delta g\,(n_{x,k}\,\hat{x} + n_{y,k}\,\hat{y} +
n_{z,k}\,\hat{z}),
\label{eq:gradient-QAM}
\end{equation}
where $n_{x,k}, n_{y,k}, n_{z,k}$ are integer lattice coordinates and
$\Delta g$ is the minimum gradient spacing.

For a cubic $M^{1/3} \times M^{1/3} \times M^{1/3}$ constellation, each
symbol carries $\log_2 M$ bits.
The minimum distance is $\Delta g$, and the BER for nearest-neighbor
decoding is
\begin{equation}
\BER_{\nabla\psi} \approx \frac{6(M^{1/3}-1)}{3M^{1/3}\log_2 M}\,
Q\!\left(\frac{\Delta g}{2\sigma_{\nabla\psi}}\right),
\label{eq:BER-grad}
\end{equation}
where $\sigma_{\nabla\psi}$ is the gradient measurement noise per axis.

\subsection{Composite $(\delta\psi, \dot{\psi}, \nabla\psi)$
Constellation}

The full 5D composite constellation jointly encodes in all degrees of
freedom.
For independent noise on each dimension, the composite BER is bounded by
the union bound:
\begin{equation}
\BER_{\rm comp} \leq \sum_{j=1}^{5} \BER_j,
\label{eq:composite-BER}
\end{equation}
where $\BER_j$ is the marginal BER on dimension $j$.

The \emph{coding gain} of the composite constellation over pure
$\delta\psi$-ASK at the same data rate is
\begin{equation}
G_{\rm cod} = 10\log_{10}\!\left(\frac{d_{\rm min,comp}^2}{d_{\rm min,ASK}^2}\right)\;\text{dB},
\label{eq:coding-gain}
\end{equation}
where $d_{\rm min}$ denotes the minimum Euclidean distance between
constellation points (after appropriate normalization).

For a 32-symbol composite constellation ($\log_2 32 = 5$ bits/symbol)
using all 5 dimensions, the coding gain over 32-ASK is approximately
\begin{equation}
G_{\rm cod}^{(32)} \approx 10\log_{10}(5) \approx 7.0\;\text{dB},
\label{eq:coding-gain-32}
\end{equation}
reflecting the sphere-packing advantage of higher-dimensional
constellations.

%=============================================================================
\section{Channel Capacity Analysis}\label{sec:capacity}
%=============================================================================

\subsection{Single-Corridor Capacity}

\begin{theorem}[Scalar $\psi$-Channel Capacity]
\label{thm:capacity}
The capacity of a single $\psi$-corridor of bandwidth $B$ with additive
Gaussian noise of spectral density $S_\psi^{\rm tot}$ and transmit
power constraint $P_\psi$ (mean-square field perturbation per unit time)
is
\begin{equation}
C_\psi = B\log_2\!\left(1 + \frac{P_\psi}{S_\psi^{\rm tot}\,B}\right)
\;\text{bits/s}.
\label{eq:capacity}
\end{equation}
\end{theorem}

\begin{proof}
The $\psi$-corridor with additive Gaussian noise is isomorphic to the
standard scalar Gaussian channel.
Shannon's coding theorem applies directly: the mutual information
$I(\delta\psi_{\rm tx};\,\delta\psi_{\rm rx})$ is maximized by Gaussian
input with power spectral density matched to the noise, yielding the
water-filling solution that reduces to~\eqref{eq:capacity} for flat
noise spectra.
\end{proof}

\paragraph{Dimensional analysis.}
The ``power'' $P_\psi$ has units of $[\psi^2/\text{s}]$, i.e., inverse
seconds (since $\psi$ is dimensionless).
The noise density $S_\psi^{\rm tot}$ has units of $[\psi^2/\text{Hz}]$.
The ratio $P_\psi / (S_\psi^{\rm tot} B)$ is the signal-to-noise ratio
per channel use.

\subsection{Multi-Dimensional Capacity}

For the full 5D composite channel $(\delta\psi, \dot\psi,
\nabla_x\psi, \nabla_y\psi, \nabla_z\psi)$ with independent noise on
each dimension:
\begin{equation}
C_{\rm comp} = \sum_{j=1}^{5} B_j\log_2\!\left(1 +
\frac{P_j}{S_j\,B_j}\right),
\label{eq:capacity-5D}
\end{equation}
where $B_j$, $P_j$, $S_j$ are the bandwidth, power, and noise density
on dimension $j$.

With equal allocation $P_j = P_\psi/5$, $B_j = B$, $S_j = S_\psi^{\rm tot}$:
\begin{equation}
C_{\rm comp} = 5B\log_2\!\left(1 + \frac{P_\psi}{5\,S_\psi^{\rm tot}\,B}\right).
\label{eq:capacity-equal}
\end{equation}
In the high-SNR regime ($\SNR \gg 1$):
\begin{equation}
C_{\rm comp} \approx 5B\log_2\!\left(\frac{\SNR}{5}\right)
\approx 5C_\psi - 5B\log_2 5.
\label{eq:capacity-highSNR}
\end{equation}
The 5D channel provides approximately $5\times$ the capacity of the
scalar channel at high SNR, with a modest penalty from power splitting.

\subsection{Water-Filling Optimization}

Optimal power allocation across the 5 dimensions follows the
water-filling solution.
Let $N_j = S_j B_j$ be the noise power on dimension $j$.
The optimal power allocation is
\begin{equation}
P_j^* = \max\!\left(\nu - N_j,\;0\right),
\label{eq:waterfilling}
\end{equation}
where the ``water level'' $\nu$ is determined by the total power constraint
$\sum_j P_j^* = P_\psi$.

In practice, the gradient dimensions $j = 3,4,5$ typically have higher
noise (gradiometry is harder than local $\psi$ measurement), so the
water-filling solution preferentially allocates power to the amplitude
and phase-rate dimensions.

\subsection{Capacity Comparison with Conventional Channels}

\begin{table}[t]
\caption{Channel capacity comparison at representative parameters.
$P_{\rm tx}$: transmit power; $B$: bandwidth; $L$: link distance;
$C$: achievable capacity. Parameters chosen for comparable link
scenarios.}
\label{tab:capacity-comparison}
\begin{ruledtabular}
\begin{tabular}{lcccc}
Modality & $P_{\rm tx}$ & $B$ & $L$ & $C$ \\
\hline
RF (5G mmW) & 1 W & 800 MHz & 200 m & 4 Gbps \\
FSO (1550 nm) & 100 mW & 10 GHz & 1 km & 40 Gbps \\
QKD (BB84) & 10 mW & --- & 100 km & 10 kbps \\
Underwater RF & 100 W & 100 kHz & 10 m & 100 kbps \\
Underwater acoustic & 10 W & 50 kHz & 1 km & 50 kbps \\
Neutrino & 1 GW & $\sim$Hz & Earth & $\sim 0.1$ bps \\
\hline
$\psi$-corridor$^\dagger$ & $P_\psi$ & $B_\psi$ & any & $5B_\psi\log_2(1+\SNR_\psi/5)$ \\
\end{tabular}
\end{ruledtabular}
\begin{flushleft}
\footnotesize{$^\dagger$ The $\psi$-corridor capacity depends on the
achievable $\psi$-pump power and detector sensitivity.
See Section~\ref{sec:demonstrator} for concrete parameter estimates.}
\end{flushleft}
\end{table}

The key advantage of $\psi$-comms is not raw capacity but rather the
combination of:
\begin{itemize}
\item Medium penetration (through rock, water, atmosphere),
\item Synchronization-free ranging (one-way delay measurement),
\item Native geometry-locked authentication,
\item Immunity to EM jamming and interception.
\end{itemize}

%=============================================================================
\section{Security Analysis: Geometry-Locked Authentication}\label{sec:security}
%=============================================================================

\subsection{Threat Model}

We consider an adversary (Eve) who:
\begin{enumerate}[label=(E\arabic*)]
\item Can observe the $\psi$-field at points outside the corridor,
\item Can generate arbitrary EM fields,
\item Can inject $\psi$-perturbations from her own $\psi$-generator,
\item \textbf{Cannot} reproduce the exact source geometry of the
legitimate transmitter (Alice).
\end{enumerate}

The key security assumption is that $\psi$ satisfies an elliptic field
equation~\eqref{eq:linearized}: the field at any point is determined by
the complete source distribution.
An adversary at a different location produces a fundamentally different
gradient map, even if the scalar amplitude can be matched at isolated
points.

\subsection{GLRT Authentication}

\begin{definition}[Corridor Template]
A corridor template $\Template_\Corridor$ is the set of $N_s$ sampled
fingerprint vectors
\begin{equation}
\Template_\Corridor = \bigl\{\Fingerprint_\Corridor(s_1),\ldots,
\Fingerprint_\Corridor(s_{N_s})\bigr\}
\label{eq:template}
\end{equation}
stored securely at the receiver during a trusted initialization phase.
\end{definition}

The receiver computes the generalized likelihood ratio test (GLRT)
statistic:
\begin{equation}
\Lambda = \sum_{i=1}^{N_s} \frac{\|\Fingerprint_{\rm meas}(s_i) -
\Fingerprint_\Corridor(s_i)\|^2}{\sigma_\Fingerprint^2(s_i)},
\label{eq:GLRT}
\end{equation}
where $\sigma_\Fingerprint^2(s_i)$ is the measurement noise variance at
sample point $s_i$.

Under the null hypothesis $\mathcal{H}_0$ (legitimate transmitter),
$\Lambda \sim \chi^2_{5N_s}$.
Under the alternative $\mathcal{H}_1$ (spoofed signal from a different
geometry), $\Lambda$ follows a non-central $\chi^2$ with non-centrality
parameter
\begin{equation}
\lambda_{\rm nc} = \sum_{i=1}^{N_s}
\frac{\|\Delta\Fingerprint_{\rm Eve}(s_i)\|^2}{\sigma_\Fingerprint^2(s_i)},
\label{eq:non-centrality}
\end{equation}
where $\Delta\Fingerprint_{\rm Eve}$ is the fingerprint deviation
produced by Eve's off-geometry source.

\begin{theorem}[Authentication Security]\label{thm:security}
For a GLRT threshold $\eta$ set to achieve false-alarm probability
$P_{\rm FA} = \alpha$, the false-acceptance probability of a spoofed
signal satisfies
\begin{equation}
P_{\rm FA,spoof} \leq \exp\!\left(-\frac{\lambda_{\rm nc}}{2}
\left(1 - \sqrt{\frac{5N_s}{\lambda_{\rm nc}}}\right)^2\right)
\label{eq:security-bound}
\end{equation}
for $\lambda_{\rm nc} > 5N_s$.
\end{theorem}

\begin{proof}
Apply the Chernoff bound to the non-central $\chi^2$ distribution.
The moment-generating function of $\chi^2_\nu(\lambda_{\rm nc})$
evaluated at $t = 1/2 - \nu/(2\lambda_{\rm nc})$ yields the
bound~\eqref{eq:security-bound} after standard optimization.
\end{proof}

\paragraph{Scaling with corridor complexity.}
The non-centrality parameter $\lambda_{\rm nc}$ grows linearly with the
number of sampling points $N_s$ (for uniformly distributed fingerprint
deviations).
The false-acceptance probability therefore decreases
\emph{exponentially} with corridor complexity:
\begin{equation}
P_{\rm FA,spoof} \sim e^{-\Theta(N_s)}.
\label{eq:exponential-security}
\end{equation}

\subsection{Impossibility of Geometric Forgery}

\begin{theorem}[Geometric Spoofing Bound]\label{thm:forgery}
Let $\delta\psi_A(\mathbf{r})$ be the field produced by Alice's source
distribution $\delta\rho_A(\mathbf{r})$ centered at $\mathbf{r}_A$, and
$\delta\psi_E(\mathbf{r})$ be Eve's field from $\delta\rho_E(\mathbf{r})$
centered at $\mathbf{r}_E \neq \mathbf{r}_A$.
Then for any $\delta\rho_E$:
\begin{equation}
\int_\Corridor \bigl|\nabla(\delta\psi_A - \delta\psi_E)\bigr|^2\,ds
\geq \frac{C_0}{|\mathbf{r}_A - \mathbf{r}_E|^2}\int_\Corridor
|\nabla\delta\psi_A|^2\,ds,
\label{eq:forgery-bound}
\end{equation}
where $C_0 > 0$ is a constant depending on the corridor geometry and
$L/|\mathbf{r}_A - \mathbf{r}_E|$.
\end{theorem}

\begin{proof}[Proof sketch]
The gradient of the Newtonian potential from a point source at
$\mathbf{r}_0$ scales as $\nabla\delta\psi \propto (\mathbf{r} -
\mathbf{r}_0)/|\mathbf{r} - \mathbf{r}_0|^3$.
For $\mathbf{r}_E \neq \mathbf{r}_A$, the angular structure of
$\nabla\delta\psi_E$ necessarily differs from $\nabla\delta\psi_A$ along
$\Corridor$.
The quadratic form on the left is minimized by the best-fit multipole
expansion, but the dipole and higher moments cannot simultaneously cancel
at all points along a one-dimensional curve, yielding the lower bound.
\end{proof}

This theorem formalizes the key security property: the spatial gradient
pattern along a corridor constitutes a physical authentication key that
cannot be forged from a different location.

\subsection{Comparison with QKD Security}

Quantum key distribution (BB84, E91, etc.) derives security from quantum
no-cloning.
$\psi$-corridor authentication derives security from the elliptic
structure of the field equation---a classical but equally fundamental
constraint.

Key differences:
\begin{itemize}
\item \textbf{Distance:} QKD is limited to $\sim 500$~km in fiber
(without quantum repeaters); $\psi$-comms have no distance limit in
principle.
\item \textbf{Medium:} QKD requires low-loss optical channels;
$\psi$-corridors penetrate all media.
\item \textbf{Rate:} QKD key rates are $\sim 10$~kbps at 100 km;
$\psi$-comms data rates depend on $\psi$-pump power.
\item \textbf{Vulnerability:} QKD is vulnerable to detector blinding and
Trojan-horse attacks on the hardware;
$\psi$-authentication is a physical-layer property immune to such side
channels.
\item \textbf{Assumption:} QKD assumes quantum mechanics;
$\psi$-authentication assumes DFD field equations.
\end{itemize}

%=============================================================================
\section{Orthogonal Corridor Multiplexing}\label{sec:multiplexing}
%=============================================================================

\subsection{Corridor Orthogonality}

\begin{definition}[Orthogonal Corridors]
Two corridors $\Corridor_1$, $\Corridor_2$ are orthogonal if their
gradient maps satisfy
\begin{equation}
\langle \nabla\delta\psi_1,\;\nabla\delta\psi_2 \rangle_\Omega
\equiv \int_\Omega \nabla\delta\psi_1 \cdot \nabla\delta\psi_2\,d^3r
= 0,
\label{eq:orthogonality}
\end{equation}
where $\Omega$ is the receiver's detection volume.
\end{definition}

Orthogonal corridors carry independent information streams without
mutual interference.
The number of orthogonal corridors accessible at a receiver is
determined by the spatial mode structure of the detection volume.

\subsection{Maximum Number of Orthogonal Corridors}

\begin{proposition}[Mode Count]
\label{prop:mode-count}
For a spherical detection volume of radius $R_d$ and corridor
wavelength $\lambda_\psi$, the number of distinguishable orthogonal
corridors arriving from different directions is
\begin{equation}
N_c \leq \left\lfloor\frac{4\pi R_d^2}{\lambda_\psi^2}\right\rfloor
\approx 4\pi\!\left(\frac{R_d}{\lambda_\psi}\right)^2.
\label{eq:mode-count}
\end{equation}
\end{proposition}

This is the angular resolution limit: corridors arriving from directions
separated by less than $\lambda_\psi / R_d$ radians are not
distinguishable by gradient measurements alone.

\subsection{$N_c$-Corridor Aggregate Capacity}

With $N_c$ orthogonal corridors and equal power allocation:
\begin{equation}
C_{\rm agg} = N_c\,C_\psi = N_c\,B\log_2\!\left(1 +
\frac{P_\psi}{N_c\,S_\psi^{\rm tot}\,B}\right).
\label{eq:aggregate-capacity}
\end{equation}
In the interference-free regime, the aggregate capacity scales linearly
with $N_c$:
\begin{equation}
C_{\rm agg} \approx N_c\,B\log_2(\SNR_\psi/N_c)
\label{eq:aggregate-highSNR}
\end{equation}
for $\SNR_\psi \gg N_c$.

\subsection{Multi-User Access: $\psi$-CDMA}

For scenarios where corridor orthogonality is imperfect, we define
$\psi$-CDMA: each user $k$ is assigned a unique gradient code
$\mathbf{g}_k(\mathbf{r})$ satisfying
\begin{equation}
\int_\Omega \mathbf{g}_j(\mathbf{r}) \cdot \mathbf{g}_k(\mathbf{r})\,d^3r
= \delta_{jk}\,\|\mathbf{g}_k\|^2.
\label{eq:psi-CDMA}
\end{equation}

The receiver decorrelates users by matched filtering against the
gradient codes.
The near-far problem is naturally mitigated because each corridor's
gradient pattern is geometrically determined by its source location,
providing automatic power-distance decorrelation.

\subsection{Stealth Corridor Hopping}

For low-probability-of-intercept (LPI) communications, the transmitter
rapidly rekeys the corridor gradient map by switching between
pre-agreed source configurations.
The hopping sequence $\{k_1, k_2, \ldots\}$ indexes a codebook of
$N_{\rm hop}$ corridor patterns:
\begin{equation}
\Corridor(t) = \Corridor_{k_{\lfloor t/T_h\rfloor}},
\label{eq:hopping}
\end{equation}
where $T_h$ is the hop duration.

The processing gain against an eavesdropper who does not know the
hopping sequence is
\begin{equation}
G_{\rm hop} = 10\log_{10}(N_{\rm hop})\;\text{dB}.
\label{eq:hopping-gain}
\end{equation}

An eavesdropper must simultaneously monitor $N_{\rm hop}$ corridor
patterns, diluting her SNR by $1/N_{\rm hop}$ per corridor.

%=============================================================================
\section{Comparison with Existing Communication Modalities}
\label{sec:comparison}
%=============================================================================

\subsection{RF and Microwave Communications}

\paragraph{Advantages of RF.}
Mature technology, high data rates (multi-Gbps in 5G/6G), standardized
protocols, broad commercial ecosystem.

\paragraph{Limitations addressed by $\psi$-comms.}
RF is vulnerable to:
\begin{itemize}
\item \textbf{Jamming:} A $\psi$-corridor is immune to EM jamming since
the information carrier is the scalar field, not the transverse EM field.
\item \textbf{Spoofing:} RF spoofing requires only matching the signal
waveform; $\psi$-spoofing requires matching the spatial gradient geometry.
\item \textbf{Multipath:} RF multipath causes fading; $\psi$-corridors
follow geodesics of the optical metric with deterministic delay.
\item \textbf{Medium penetration:} RF penetrates poorly through water
($\sim$meters at UHF), rock ($\sim$tens of meters), and metal (not at
all); $\psi$ penetrates all media.
\end{itemize}

Research on beyond-5G/6G systems is exploring terahertz frequencies
(0.1--10 THz), orbital angular momentum multiplexing, and reconfigurable
intelligent surfaces.
These advances increase capacity but do not address the fundamental
vulnerabilities of EM-carrier communications.

\subsection{Free-Space Optical Communications}

Free-space optical (FSO) links offer very high data rates (100+ Gbps)
but require line-of-sight, are degraded by atmospheric turbulence, fog,
and precipitation, and provide no native authentication.
$\psi$-corridors are unaffected by atmospheric conditions and provide
geometry-locked authentication inherently.

\subsection{Quantum Key Distribution}

QKD provides information-theoretic security for key distribution but
faces fundamental practical challenges:
\begin{itemize}
\item \textbf{Distance:} Limited to $\sim 500$ km in fiber without
quantum repeaters (which remain experimental).
\item \textbf{Rate:} Key generation rates of $\sim 10$ kbps at
100 km, scaling inversely with distance.
\item \textbf{Detector vulnerabilities:} Blinding attacks, after-gate
attacks, and Trojan-horse attacks exploit detector imperfections.
\item \textbf{Side channels:} Practical QKD systems leak information
through timing, power consumption, and EM emissions.
\item \textbf{Medium:} Requires low-loss optical paths (fiber or
free-space line-of-sight).
\end{itemize}

$\psi$-corridor authentication provides physical-layer security that
is complementary to QKD: while QKD secures the key exchange,
$\psi$-authentication secures the link identity.
A hybrid QKD/$\psi$-corridor system could provide both key security and
link authentication.

\subsection{Underwater Communications}

Underwater communication is severely constrained:
\begin{itemize}
\item \textbf{EM:} Seawater attenuates RF rapidly; practical range is
$\sim 10$ m at UHF.
\item \textbf{Acoustic:} Limited bandwidth ($\sim 50$ kHz), severe
multipath, Doppler sensitivity, range $\sim 1$--$10$ km.
\item \textbf{Optical:} Blue-green lasers penetrate to $\sim 200$ m in
clear water; severely limited in turbid conditions.
\end{itemize}

$\psi$-corridors penetrate seawater with zero electromagnetic attenuation.
The scalar field couples to mass-energy density, not EM properties, so
the salt content and turbidity of water are irrelevant.
This makes $\psi$-comms uniquely suited for submarine communication,
underwater sensor networks, and deep-ocean exploration.

\subsection{Deep-Space Communications}

Current deep-space links (e.g., NASA Deep Space Network) use X-band and
Ka-band RF, achieving data rates of $\sim 1$--$10$ Mbps at Mars
distance and $\sim 1$--$100$ kbps at outer planet distances.
The primary limitations are:
\begin{itemize}
\item $1/r^2$ free-space path loss,
\item Solar conjunction blackout periods,
\item Light-travel delay (4--20 minutes to Mars; 1--1.5 hours to Saturn).
\end{itemize}

$\psi$-corridors experience the same $1/r^2$ geometric spreading
(the field equation~\eqref{eq:linearized} has the same Green's function
as the Coulomb potential), but offer:
\begin{itemize}
\item Penetration through the solar corona (no conjunction blackout),
\item One-way ranging without two-way light-time exchange,
\item Native authentication preventing signal spoofing during critical
operations (e.g., autonomous spacecraft commanding).
\end{itemize}

\subsection{Neutrino and Gravitational-Wave Communications}

Neutrino communication has been experimentally demonstrated at extremely
low data rates ($\sim 0.1$ bits/s) using the NuMI neutrino beam at
Fermilab and the MINERvA detector.
The system required a multi-megawatt accelerator and a 170-ton
detector---clearly impractical for deployment.

Gravitational-wave communication concepts have been proposed but require
generating detectable spacetime curvature perturbations, which demands
astrophysical-scale energy ($\sim M_\odot c^2$).

$\psi$-corridor communication occupies a unique intermediate niche:
the coupling is via the scalar sector (which may be excited at laboratory
scales if $|\lambda - 1| \neq 0$), and the detection exploits the
same high-precision metrology developed for gravitational-wave detectors
but applied to the scalar mode.

\begin{table*}[t]
\caption{Comprehensive comparison of communication modalities.
Categories: capacity, security, medium penetration, and deployment
readiness.}
\label{tab:comparison}
\begin{ruledtabular}
\begin{tabular}{lcccccc}
Property & RF/MW & FSO & QKD & Underwater Ac. & Neutrino &
$\psi$-Corridor \\
\hline
Max data rate & $\sim$100 Gbps & $\sim$100 Gbps & $\sim$10 kbps &
$\sim$50 kbps & $\sim$0.1 bps & TBD$^a$ \\
Medium penetration & Poor & None & None & Moderate & Total & Total \\
EM jam immunity & No & No & Partial$^b$ & Yes & Yes & Yes \\
Native authentication & No & No & QM-based & No & No & Geometry-locked \\
Anti-spoofing & Weak & Weak & Strong & Weak & N/A & Strong \\
Range limit & LOS/reflect & LOS & $\sim$500 km & $\sim$10 km &
$\sim$Earth & None in principle \\
Synch-free ranging & No & No & No & No & No & Yes \\
Tech readiness & TRL 9 & TRL 7 & TRL 5 & TRL 7 & TRL 2 & TRL 1 \\
\end{tabular}
\end{ruledtabular}
\begin{flushleft}
\footnotesize{$^a$ Depends on achievable $\psi$-pump power and detector
sensitivity; see Section~\ref{sec:demonstrator}.\\
$^b$ QKD itself is not jammed by EM, but it requires optical channels
vulnerable to physical disruption.}
\end{flushleft}
\end{table*}

%=============================================================================
\section{Application-Specific Designs}\label{sec:applications}
%=============================================================================

\subsection{Secure Military Communications}

\subsubsection{Requirements}

\begin{itemize}
\item Anti-jam, anti-spoof, low probability of intercept (LPI),
\item Operation in GPS-denied environments,
\item Penetration through urban structures, underground, and underwater.
\end{itemize}

\subsubsection{Architecture}

A military $\psi$-comm network consists of:
\begin{enumerate}
\item \textbf{Command nodes:} High-power $\psi$-generators producing
steered corridors to subordinate units.
\item \textbf{Field terminals:} Compact $\psi$-detectors with
gradiometry capability for corridor identification and demodulation.
\item \textbf{Network controller:} Manages corridor allocation, hopping
sequences, and route stitching.
\end{enumerate}

Stealth corridor hopping (Section~\ref{sec:multiplexing}) provides LPI
with processing gain $G_{\rm hop} = 10\log_{10}(N_{\rm hop})$ dB.
For $N_{\rm hop} = 1000$: $G_{\rm hop} = 30$ dB.

Geometry-locked authentication replaces cryptographic IFF
(identification friend-or-foe), providing physical-layer authentication
that cannot be compromised by captured encryption keys.

\subsection{Financial Timing Backbone}

\subsubsection{Requirements}

\begin{itemize}
\item Sub-nanosecond time synchronization across continental distances,
\item Spoof-proof timing (preventing time-spoofing attacks on
high-frequency trading),
\item Independence from GPS/GNSS (vulnerability to jamming and spoofing).
\end{itemize}

\subsubsection{Architecture}

The $\psi$-corridor provides one-way timing via the intrinsic delay:
\begin{equation}
\tau = \frac{1}{c}\int e^\psi\,ds.
\end{equation}
Two key advantages:
\begin{enumerate}
\item \textbf{One-way ranging:} No two-way exchange needed; the delay
is determined by the known corridor geometry and the fundamental
constant $c$.
\item \textbf{Anti-spoofing:} A timing spoofer must reproduce the exact
$\psi$-gradient profile, not merely the time-of-arrival.
\end{enumerate}

For a financial timing backbone linking New York and Chicago
($L \approx 1200$ km):
\begin{equation}
\tau_0 = \frac{L}{c}(1 + \psi_0) \approx 4.0\;\text{ms},
\end{equation}
where $\psi_0 \sim 10^{-9}$ (Earth's surface potential).
The timing precision is limited by the $\psi$-measurement noise, not by
propagation medium variability.

\subsection{Underwater Sensor Networks}

\subsubsection{Architecture}

An underwater $\psi$-network consists of:
\begin{enumerate}
\item \textbf{Surface gateway:} High-power $\psi$-generator at the
ocean surface.
\item \textbf{Subsea nodes:} Autonomous underwater vehicles (AUVs) and
fixed sensors with $\psi$-detectors.
\item \textbf{Acoustic hybrid:} Optional acoustic payload carried within
the $\psi$-corridor for high-bandwidth data.
\end{enumerate}

The $\psi$-corridor provides:
\begin{itemize}
\item Reliable command/control through the water column (independent of
depth, salinity, temperature),
\item Synchronization-free localization (each node measures its
$\psi$-corridor delay to multiple surface gateways),
\item Deterministic multipath: the $\psi$-corridor geometry is set by
the source distribution, not by acoustic scattering.
\end{itemize}

\subsection{Space-to-Ground Links}

\subsubsection{Orbital Corridor Architecture}

Satellites equipped with $\psi$-generators create orbital corridor meshes
to ground stations.
Applications include:
\begin{itemize}
\item \textbf{GNSS augmentation:} $\psi$-corridor timing overlays
providing spoof-proof position and timing.
\item \textbf{Deep-space relay:} $\psi$-corridor links bypassing
solar conjunction for continuous Mars/Jupiter communication.
\item \textbf{Satellite commanding:} Geometry-locked authentication
preventing unauthorized spacecraft commanding.
\end{itemize}

The one-way delay from a satellite at altitude $h$:
\begin{equation}
\tau_{\rm sat} = \frac{1}{c}\int_0^h e^{\psi(r)}\,dr
\approx \frac{h}{c}\left(1 + \frac{GM}{c^2 R_\oplus}\ln\frac{R_\oplus + h}{R_\oplus}\right),
\label{eq:sat-delay}
\end{equation}
where the logarithmic correction encodes the gravitational $\psi$ profile.
For LEO ($h = 400$ km): $\tau_{\rm sat} \approx 1.33$ ms with a
gravitational correction of $\sim 50$ ps---detectable with current
optical clock technology.

\subsection{Emergency Mesh Communications}

\subsubsection{Architecture}

After catastrophic infrastructure loss (earthquake, EMP, nuclear event),
$\psi$-mesh nodes form ad hoc corridors:
\begin{enumerate}
\item Each node activates its $\psi$-generator, creating corridors to
all detectable neighbors.
\item Distributed scheduling allocates orthogonal corridor slots.
\item Routes are stitched hop-by-hop with corridor fingerprint
authentication at each relay.
\end{enumerate}

The key advantage over RF mesh networks: $\psi$-corridors penetrate
through rubble, collapsed buildings, and flooded tunnels, enabling
communication in scenarios where all RF paths are blocked.

%=============================================================================
\section{Laboratory Demonstrator Design}\label{sec:demonstrator}
%=============================================================================

\subsection{Design Philosophy}

The demonstrator aims to establish proof-of-principle for three core
capabilities:
\begin{enumerate}
\item $\psi$-field generation via EM--$\psi$ pumping,
\item Information encoding in $\delta\psi$ modulation,
\item Corridor fingerprint measurement for authentication.
\end{enumerate}

The design leverages the EM--$\psi$ coupling framework of Appendix R
in~\cite{Alcock2025DFD}, specifically the parametric amplification
channel at $2\omega = 2\Omega_\psi$.

\subsection{System Parameters}

\begin{table}[h]
\caption{Point-to-point demonstrator parameters.}
\label{tab:demonstrator}
\begin{ruledtabular}
\begin{tabular}{lcc}
Parameter & Symbol & Value \\
\hline
Cavity stored energy & $U_0$ & $10^5$ J (100 kJ) \\
Cavity $Q$-factor & $Q$ & $10^{10}$ \\
Cavity frequency & $\omega/2\pi$ & 1.3 GHz \\
Modulation depth & $m$ & 0.1 \\
$\psi$-mode frequency & $\Omega_\psi/2\pi$ & $\sim 10$ Hz \\
$\psi$-mode loss & $\gamma_\psi/\Omega_\psi$ & $10^{-3}$ \\
Tube cross-section & $A_\psi$ & 0.27 m$^2$ \\
Cavity aperture total & $A_{\rm cav,tot}$ & $3 \times 10^{-2}$ m$^2$ \\
Link distance & $L$ & 10 m \\
Detector type & --- & SQUID gradiometer \\
Detector sensitivity & $\sigma_{\nabla\psi}$ & $10^{-14}$ m$^{-1}$Hz$^{-1/2}$ \\
\end{tabular}
\end{ruledtabular}
\end{table}

\subsection{Transmitter: $\psi$-Pump}

The transmitter consists of an SRF (superconducting radio-frequency)
cavity array operated in TE+TM superposition mode to break the geometry
transparency condition and restore the driven overlap $G \neq 0$.

\subsubsection{Parametric $\psi$-Generation}

The Mathieu gain parameter from Eq.~(R.6) of~\cite{Alcock2025DFD}:
\begin{equation}
h = (\lambda - 1)\frac{U_0}{M_\psi \Omega_\psi^2}Hm.
\label{eq:Mathieu-gain}
\end{equation}

For the demonstrator parameters (Table~\ref{tab:demonstrator}), with
$M_\psi \approx A_\psi L_{\rm mode} / (2\pi c_s)$ and
$H \approx (2/L_{\rm mode})\kappa_{\rm eff}(A_{\rm cav,tot}/A_\psi)$:
\begin{equation}
h \approx (\lambda - 1) \times 3 \times 10^{13}.
\label{eq:h-value}
\end{equation}
For $|\lambda - 1| = 10^{-5}$ (within the accidental bound):
$h \approx 3 \times 10^8$, giving exponential parametric growth above
the damping threshold.

\subsubsection{Modulation}

Information is encoded by modulating the $\psi$-pump drive:
\begin{itemize}
\item \textbf{$\delta\psi$-ASK:} Varying the stored energy $U_0(t)$
modulates $|q(t)|$, producing amplitude variations in $\delta\psi$.
\item \textbf{$\dot{\psi}$-FSK:} Chirping the modulation frequency
around $2\Omega_\psi$ shifts the $\psi$-mode phase rate.
\item \textbf{$\nabla\psi$-key:} Switching between different cavity
array configurations changes the gradient pattern.
\end{itemize}

The maximum symbol rate is limited by the $\psi$-mode response time:
\begin{equation}
R_s \leq \gamma_\psi \sim 10^{-2}\;\text{Hz},
\label{eq:symbol-rate}
\end{equation}
yielding an initial data rate of order $10^{-2}$ bits/s for binary ASK.
While extremely low, this represents a \emph{proof-of-principle};
higher data rates require higher-frequency $\psi$ modes.

\subsection{Receiver: $\psi$-Gradiometer}

The receiver employs a SQUID-based gradiometer measuring the spatial
derivatives $\partial_i\psi$.
The $\psi$ field couples to the SQUID through its effect on local clock
rates---a SQUID loop experiences a $\psi$-dependent phase shift:
\begin{equation}
\Delta\phi_{\rm SQUID} = \frac{2eV}{\hbar}\,\delta\tau
= \frac{2eV}{\hbar}\,\frac{\delta\psi\,L_{\rm loop}}{c},
\label{eq:SQUID-phase}
\end{equation}
where $V$ is the junction voltage and $L_{\rm loop}$ is the loop
dimension.

For $V = 10\;\mu$V, $L_{\rm loop} = 0.1$ m, and $\delta\psi = 10^{-12}$:
\begin{equation}
\Delta\phi_{\rm SQUID} \approx \frac{2(1.6\times 10^{-19})(10^{-5})}
{(1.05\times 10^{-34})} \times \frac{10^{-12} \times 0.1}{3\times 10^8}
\approx 10^{-14}\;\text{rad}.
\label{eq:SQUID-sensitivity}
\end{equation}
This is at the edge of current SQUID sensitivity
($\sim 10^{-6}\;\text{rad}/\sqrt{\text{Hz}}$), requiring integration
times of $\sim 10^{16}$ s---clearly impractical with single-loop
SQUIDs.

A more practical approach uses differential atomic clock comparisons.
Two optical clocks separated by distance $d$ measure the differential
redshift:
\begin{equation}
\frac{\Delta\nu}{\nu} = \frac{\delta\psi \cdot d}{2}\,\frac{|\nabla\psi|}{d}
= \frac{|\nabla\psi| \cdot d}{2}.
\label{eq:clock-gradient}
\end{equation}
For $d = 1$ m and $|\nabla\psi| = 10^{-14}$ m$^{-1}$:
$\Delta\nu/\nu \sim 5 \times 10^{-15}$, detectable with state-of-the-art
optical clock comparisons at the $10^{-18}$ level in $\sim 100$ s
integration.

\subsection{Alternative Detector: Atom Interferometer}

A long-baseline atom interferometer (MAGIS-100 class) provides
differential $\psi$-sensing over a 100 m baseline:
\begin{equation}
\Delta\phi_{\rm AI} = k_{\rm eff}\,\frac{c^2}{2}\,
\nabla\psi \cdot L \cdot T^2,
\label{eq:AI-sensitivity}
\end{equation}
where $k_{\rm eff}$ is the effective wavevector, $L$ is the baseline,
and $T$ is the interrogation time.

For $k_{\rm eff} = 10^7$ m$^{-1}$, $L = 100$ m, $T = 1$ s,
$|\nabla\psi| = 10^{-14}$ m$^{-1}$:
\begin{equation}
\Delta\phi_{\rm AI} \sim 10^7 \times \frac{(3\times 10^8)^2}{2} \times
10^{-14} \times 100 \times 1 \sim 5 \times 10^{3}\;\text{rad}.
\label{eq:AI-value}
\end{equation}
This enormous phase shift indicates that atom interferometers are
extremely sensitive $\psi$-detectors---the challenge is
\emph{generating} the $\psi$ signal, not detecting it.

\subsection{Link Budget}

For the 10 m demonstration link:
\begin{align}
\text{Transmit field:} &\quad \delta\psi_{\rm tx} = (\lambda - 1)\,\eta\,\frac{U_0}{M_\psi \Omega_\psi \gamma_\psi} \notag \\
\text{Path loss:} &\quad \delta\psi(r) = \delta\psi_{\rm tx}\frac{R_{\rm src}^2}{r^2} \label{eq:path-loss} \\
\text{Receiver field:} &\quad \delta\psi_{\rm rx} = \delta\psi_{\rm tx}\frac{R_{\rm src}^2}{L^2} \notag
\end{align}

For $|\lambda - 1| = 10^{-5}$, $\eta = 0.3$ (TE+TM overlap),
$R_{\rm src} = 0.3$ m:
\begin{equation}
\delta\psi_{\rm rx} \sim 10^{-5} \times 0.3 \times
\frac{10^5}{M_\psi \times 60 \times 0.06} \times \frac{0.09}{100}
\sim 10^{-12},
\label{eq:rx-field}
\end{equation}
which is detectable by atomic clock gradiometry in $\sim 100$ s
integration at $10^{-18}$ fractional frequency precision.

%=============================================================================
\section{$\psi$-Field Communication Protocols}\label{sec:protocols}
%=============================================================================

\subsection{Link Initialization}

\begin{enumerate}
\item \textbf{Corridor establishment:} Transmitter activates
$\psi$-generator, sculpting corridor to receiver.
\item \textbf{Template acquisition:} Receiver measures corridor
fingerprint $\Fingerprint_\Corridor$ over $N_s$ sampling points.
\item \textbf{Template verification:} Fingerprint is verified against
predicted geometry (if source location is known) or established through
a trusted side channel.
\item \textbf{Data mode:} Transmitter begins modulating $\delta\psi$
with information content; receiver demodulates while continuously
verifying fingerprint.
\end{enumerate}

\subsection{Error Correction}

The low data rate of initial $\psi$-comms makes error correction
essential.
We propose a concatenated coding scheme:
\begin{enumerate}
\item \textbf{Inner code:} Rate-1/3 turbo code providing near-Shannon
performance at $E_b/N_0 \approx 0.5$ dB.
\item \textbf{Outer code:} Reed-Solomon RS(255,223) for burst error
protection.
\end{enumerate}

The effective coding gain is approximately 10 dB, reducing the required
receive SNR from $\sim 10$ dB (uncoded BER = $10^{-5}$) to $\sim 0$ dB.

\subsection{Timing Recovery}

Unlike RF communications, $\psi$-comms do not require two-way time
transfer for synchronization.
The one-way delay $\tau = (1/c)\int e^\psi\,ds$ is computable from the
known corridor geometry and serves as the timing reference.

Clock offset between transmitter and receiver is resolved by:
\begin{enumerate}
\item Measuring the $\psi$-corridor delay $\tau$ with the local clock.
\item Computing the expected delay from the corridor geometry.
\item The difference is the clock offset: $\Delta t = \tau_{\rm meas}
- \tau_{\rm comp}$.
\end{enumerate}

This provides one-way synchronization to the precision of the
$\psi$-measurement, which for atomic-clock detectors is
$\sigma_\tau \sim 10^{-18}$ s $\times c / |\nabla\psi| \sim
10^{-4}$ s in initial implementations.

%=============================================================================
\section{Theoretical Foundations}\label{sec:theory}
%=============================================================================

\subsection{$\psi$-Field Propagation Speed}

Perturbations $\delta\psi$ propagate according to the full dynamic
field equation including the temporal kinetic term:
\begin{equation}
\nabla^2\delta\psi - \frac{1}{c_\psi^2}\ddot{\delta\psi} =
-\frac{8\pi G}{c^2}\delta\rho,
\label{eq:wave-eq}
\end{equation}
where $c_\psi$ is the $\psi$-mode propagation speed.
In the DFD framework, $c_\psi \leq c$ with equality in the linear
regime.

The dispersion relation for plane-wave perturbations
$\delta\psi \propto e^{i(\mathbf{k}\cdot\mathbf{r} - \omega t)}$ is
\begin{equation}
\omega^2 = c_\psi^2 k^2,
\label{eq:dispersion}
\end{equation}
yielding non-dispersive propagation in the linear regime.
This is essential for communication: symbols propagate without
distortion.

\subsection{Channel Bandwidth}

The effective bandwidth of a $\psi$-corridor is set by the highest
$\psi$-mode frequency that can be excited by the transmitter and
resolved by the receiver:
\begin{equation}
B_\psi = \min(\Omega_{\psi,\max},\;\Omega_{\rm det,max}),
\label{eq:bandwidth}
\end{equation}
where $\Omega_{\psi,\max}$ is limited by the $\psi$-generator's
modulation speed and $\Omega_{\rm det,max}$ by the detector's response
time.

For the parametric pumping scheme, the achievable bandwidth is:
\begin{equation}
B_\psi \sim \gamma_\psi \sim \Omega_\psi / Q_\psi,
\label{eq:bandwidth-estimate}
\end{equation}
where $Q_\psi = \Omega_\psi / (2\gamma_\psi)$ is the $\psi$-mode
quality factor.

\subsection{Energy Requirements}

The energy required to establish a $\psi$-corridor of amplitude
$\delta\psi$ over length $L$ and cross-section $A_\psi$ is
\begin{equation}
E_\psi = \frac{\astar^2}{8\pi G}\int_\Corridor
|\nabla\delta\psi|^2\,d^3r
\approx \frac{\astar^2}{8\pi G}\frac{(\delta\psi)^2}{L} A_\psi L
= \frac{\astar^2 A_\psi (\delta\psi)^2}{8\pi G}.
\label{eq:corridor-energy}
\end{equation}

For $\astar \approx 2.7 \times 10^{-27}$ m$^{-1}$,
$A_\psi = 0.27$ m$^2$, $\delta\psi = 10^{-12}$:
\begin{equation}
E_\psi \approx \frac{(2.7\times 10^{-27})^2 \times 0.27 \times 10^{-24}}
{8\pi \times 6.67 \times 10^{-11}}
\approx 10^{-97}\;\text{J}.
\label{eq:energy-value}
\end{equation}
This extraordinarily small energy reflects the weakness of gravitational
coupling at laboratory scales.
The challenge is not the energy in the $\psi$-field itself, but the
coupling efficiency from the EM pump to the $\psi$-mode, controlled by
$|\lambda - 1|$.

The EM energy required in the pump cavity to produce $\delta\psi$ is:
\begin{equation}
U_{\rm EM} \sim \frac{E_\psi}{|\lambda - 1|\,\eta\,m} \sim
\frac{10^{-97}}{10^{-5} \times 0.3 \times 0.1} \sim 10^{-92}\;\text{J},
\label{eq:EM-energy}
\end{equation}
which is negligible---the practical limit is not energy but \emph{field
sensitivity} at the receiver.

%=============================================================================
\section{Advanced Topics}\label{sec:advanced}
%=============================================================================

\subsection{Hybrid EM/$\psi$ Links}

A hybrid link carries EM payload (high data rate) within a $\psi$-corridor
(timing, authentication, routing):
\begin{itemize}
\item The $\psi$-corridor provides the ``pipe''---timing, route identity,
and authentication.
\item The EM payload (optical or RF) provides the data bandwidth.
\item The EM carrier rides within the refractive-index structure of the
corridor.
\end{itemize}

The EM signal experiences the corridor's refractive index
$n = e^\psi$, producing a small but measurable frequency shift:
\begin{equation}
\frac{\Delta\nu_{\rm EM}}{\nu_{\rm EM}} = -\delta\psi.
\label{eq:EM-shift}
\end{equation}
This shift serves as a passive verification that the EM signal is
traveling within the intended corridor.

\subsection{Quantum Payload Tunneling}

Entangled photon pairs can be transmitted within a $\psi$-corridor,
with the corridor providing:
\begin{itemize}
\item Deterministic path (no multipath-induced decoherence),
\item Stable refractive index (reduced phase noise),
\item Geometry-locked authentication (preventing man-in-the-middle
attacks on the quantum channel).
\end{itemize}

The corridor's refractive index variation $\delta\psi \ll 1$ introduces
a phase shift $\delta\phi = \omega\,\delta\psi\,L/c$ on the entangled
photons.
For $\omega = 2\pi \times 3 \times 10^{14}$ Hz (1 $\mu$m),
$\delta\psi = 10^{-12}$, $L = 100$ km:
\begin{equation}
\delta\phi \approx 2\pi \times 3 \times 10^{14} \times 10^{-12}
\times \frac{10^5}{3\times 10^8} \approx 0.6\;\text{rad}.
\label{eq:quantum-phase}
\end{equation}
This known, deterministic phase shift can be pre-compensated, preserving
entanglement coherence.

\subsection{Corridor Network Topology}

A $\psi$-network stitches corridor segments end-to-end through relay
nodes:
\begin{equation}
\Corridor_{A \to C} = \Corridor_{A \to B} \oplus \Corridor_{B \to C},
\label{eq:route-stitch}
\end{equation}
where $\oplus$ denotes corridor concatenation with fingerprint
handoff at node $B$.

The routing metric is the total corridor delay:
\begin{equation}
\tau_{A \to C} = \tau_{A \to B} + \tau_{B \to C} +
\tau_{\rm relay}(B),
\label{eq:route-delay}
\end{equation}
where $\tau_{\rm relay}$ accounts for relay processing.

Shortest-delay routing in a $\psi$-network is isomorphic to
shortest-path routing in a weighted graph, with weights equal to
corridor delays.
Standard algorithms (Dijkstra, Bellman-Ford) apply directly.

%=============================================================================
\section{Experimental Roadmap}\label{sec:roadmap}
%=============================================================================

\subsection{Phase 1: $\psi$-Generation Verification (Years 1--3)}

Objective: Demonstrate controlled $\psi$-field generation via EM--$\psi$
parametric pumping.

\begin{itemize}
\item Deploy SRF cavity array in TE+TM superposition mode.
\item Apply deliberate $2\omega$ amplitude modulation.
\item Monitor for parametric growth at $\Omega_\psi$ using a proximate
atomic clock pair.
\item Decision: If $|\lambda - 1| \neq 0$ detected at
$> 5\sigma$, proceed to Phase 2.
\end{itemize}

\subsection{Phase 2: Point-to-Point Demonstration (Years 3--5)}

Objective: Transmit and receive information over a 10 m $\psi$-corridor.

\begin{itemize}
\item Modulate $\psi$-pump with binary ASK at $\sim 0.01$ Hz.
\item Detect at 10 m with differential atomic clock pair.
\item Demonstrate corridor fingerprint measurement.
\item Metric: BER $< 10^{-3}$ at any data rate.
\end{itemize}

\subsection{Phase 3: Authentication Demonstration (Years 5--7)}

Objective: Demonstrate geometry-locked authentication.

\begin{itemize}
\item Establish corridor template during trusted initialization.
\item Attempt spoofing from an offset transmitter location.
\item Measure GLRT rejection of spoofed signals.
\item Metric: False-acceptance rate $< 10^{-6}$ for
$|\Delta\mathbf{r}_{\rm spoof}| > 1$ m.
\end{itemize}

\subsection{Phase 4: Underwater and Through-Medium (Years 7--10)}

Objective: Demonstrate $\psi$-corridor communication through opaque media.

\begin{itemize}
\item Deploy transmitter and receiver separated by water tank, rock
barrier, or Faraday cage.
\item Verify zero EM leakage (Faraday cage test).
\item Demonstrate data transmission through medium.
\end{itemize}

%=============================================================================
\section{Conclusions}\label{sec:conclusions}
%=============================================================================

We have developed a complete communication theory for $\psi$-field
information transfer via scalar refractive corridors, as specified in
Patent Application~\cite{Alcock2025Patent8}.
The key results are:

\begin{enumerate}
\item \textbf{Channel capacity:} The 5D composite
$(\delta\psi, \dot\psi, \nabla\psi)$ constellation achieves
approximately $5\times$ the capacity of scalar modulation alone,
with the Shannon limit~\eqref{eq:capacity-5D} bounding achievable
rates.

\item \textbf{Security:} Geometry-locked authentication via corridor
fingerprint GLRT matching achieves exponentially vanishing
false-acceptance probability
$P_{\rm FA,spoof} \sim e^{-\Theta(N_s)}$ with corridor complexity
(Theorem~\ref{thm:security}).
The impossibility of geometric forgery
(Theorem~\ref{thm:forgery}) provides a classical security guarantee
analogous to quantum no-cloning.

\item \textbf{Multiplexing:} Orthogonal corridor multiplexing supports
up to $N_c \sim 4\pi(R_d/\lambda_\psi)^2$ independent channels,
with $\psi$-CDMA providing multi-user access when orthogonality is
imperfect.

\item \textbf{Unique niche:} $\psi$-communications occupy a
communication space inaccessible to all existing modalities:
medium-penetrating, synchronization-free, natively authenticated
links immune to EM jamming (Table~\ref{tab:comparison}).

\item \textbf{Demonstrator:} A point-to-point link at 10 m using
SRF $\psi$-pumping and atomic clock detection is feasible at the
current EM--$\psi$ coupling bound
$|\lambda - 1| \lesssim 3 \times 10^{-5}$.
\end{enumerate}

The critical prerequisite is experimental verification that
$|\lambda - 1| \neq 0$---that electromagnetic fields can actively
source $\psi$-mode oscillations.
If confirmed, $\psi$-field communication would represent a
fundamentally new physical layer for information transfer, with
applications spanning secure military communications, financial
timing, underwater networks, deep-space links, and disaster-resilient
emergency mesh systems.

%=============================================================================
% REFERENCES
%=============================================================================

\begin{thebibliography}{99}

\bibitem{Alcock2025DFD}
G.~T.~Alcock,
``Density Field Dynamics: A Complete Unified Theory,'' v3.2 (2025).

\bibitem{Alcock2025Patent8}
G.~T.~Alcock,
``$\psi$-Field Communication and Information Systems,''
U.S. Provisional Patent Application (2025).

\bibitem{Gordon1923}
W.~Gordon,
``Zur Lichtfortpflanzung nach der Relativit\"atstheorie,''
\textit{Ann. Phys.} \textbf{72}, 421 (1923).

\bibitem{Perlick2000}
V.~Perlick,
\textit{Ray Optics, Fermat's Principle, and Applications to General Relativity}
(Springer, 2000).

\bibitem{Shannon1948}
C.~E.~Shannon,
``A mathematical theory of communication,''
\textit{Bell Syst. Tech. J.} \textbf{27}, 379 (1948).

\bibitem{BekensteinMilgrom1984}
J.~D.~Bekenstein and M.~Milgrom,
``Does the missing mass problem signal the breakdown of Newtonian gravity?''
\textit{Astrophys. J.} \textbf{286}, 7 (1984).

\bibitem{McGaugh2016RAR}
S.~S.~McGaugh, F.~Lelli, and J.~M.~Schombert,
``Radial Acceleration Relation in Rotationally Supported Galaxies,''
\textit{Phys. Rev. Lett.} \textbf{117}, 201101 (2016).

\bibitem{NuMI2012}
S.~Stancil \textit{et al.},
``Demonstration of Communication Using Neutrinos,''
\textit{Mod. Phys. Lett. A} \textbf{27}, 1250077 (2012).

\bibitem{GW170817}
B.~P.~Abbott \textit{et al.} (LIGO/Virgo),
``Gravitational Waves and Gamma-Rays from a Binary Neutron Star Merger: GW170817 and GRB 170817A,''
\textit{Astrophys. J. Lett.} \textbf{848}, L13 (2017).

\bibitem{BB84}
C.~H.~Bennett and G.~Brassard,
``Quantum cryptography: Public key distribution and coin tossing,''
\textit{Theor. Comput. Sci.} \textbf{560}, 7 (2014).

\bibitem{ScaraniQKD}
V.~Scarani \textit{et al.},
``The security of practical quantum key distribution,''
\textit{Rev. Mod. Phys.} \textbf{81}, 1301 (2009).

\bibitem{MAGIS100}
M.~Abe \textit{et al.},
``Matter-wave Atomic Gradiometer Interferometric Sensor (MAGIS-100),''
\textit{Quantum Sci. Technol.} \textbf{6}, 044003 (2021).

\bibitem{DSN}
J.~Taylor, D.~K.~Lee, and S.~Shambayati,
``Mars Reconnaissance Orbiter Telecommunications,''
\textit{DESCANSO Design and Performance Summary Series} (2006).

\bibitem{UnderwaterComm}
M.~Chitre, S.~Shahabudeen, and M.~Stojanovic,
``Underwater acoustic communications and networking: Recent advances and future challenges,''
\textit{Marine Technology Society Journal} \textbf{42}, 103 (2008).

\bibitem{6GVision}
W.~Saad, M.~Bennis, and M.~Chen,
``A Vision of 6G Wireless Systems: Applications, Trends, Technologies, and Open Research Problems,''
\textit{IEEE Network} \textbf{34}, 134 (2020).

\bibitem{TerahertzComm}
T.~K\"urner and S.~Priebe,
``Towards THz Communications -- Status in Research, Standardization and Regulation,''
\textit{J. Infrared Millim. THz Waves} \textbf{35}, 53 (2014).

\bibitem{OAM}
J.~Wang \textit{et al.},
``Terabit free-space data transmission employing orbital angular momentum multiplexing,''
\textit{Nature Photon.} \textbf{6}, 488 (2012).

\bibitem{DetectorBlinding}
L.~Lydersen \textit{et al.},
``Hacking commercial quantum cryptography systems by tailored bright illumination,''
\textit{Nature Photon.} \textbf{4}, 686 (2010).

\end{thebibliography}

\end{document}
